% !TEX root = ../../main/main.tex
\hypeexplica{../sections/hype-explica-erro-tipo-i-ii/banner-erro-tipo-i-ii.png}
{Gustavo Neves}
{

\hspace*{1.5em} Todo teste estatístico está sujeito a erro,
isso é inevitável quando se trabalha com amostra em vez de população inteira.
O que poucas pessoas param pra notar é que existem dois tipos bem diferentes de erro
escondidos atrás de qualquer "p < 0,05",
e confundir um com o outro muda completamente a conclusão de um estudo,
de um teste A/B ou de um exame médico.

\subtitle{Relembrando o teste de hipótese}

\hspace*{1.5em} Todo teste de hipótese parte de uma pergunta prática do
tipo "esse remédio funciona?" ou "essa mudança no site aumentou a conversão?".
Para responder isso de forma rigorosa,
a estatística formaliza a pergunta em duas hipóteses concorrentes, partindo da hipótese nula ($H_0$)
que normalmente representa o cenário onde "nada mudou", a diferença observada é meramente acaso,
e, por sua vez, a hipótese alternativa ($H_1$),
o efeito que se quer demonstrar que existe.

\hspace*{1.5em} Com base nos dados coletados,
o teste decide entre
rejeitar $H_0$ (concluir que há evidência de efeito)
ou não rejeitar $H_0$ (não há evidência suficiente para afirmar que existe efeito).
Vale evidenciar que a segunda opção é "não rejeitar", diferente de "aceitar como verdadeira".

\hspace*{1.5em} O ponto central é que essa decisão é sempre tomada a partir de uma amostra,
não da população inteira.
Por isso, por mais cuidadoso que o teste seja,
sempre existe uma margem real de erro.

\subtitle{Os dois jeitos de errar}

\hspace*{1.5em} Existem exatamente duas formas de tomar a decisão errada:

\begin{hypeitemize}
\item \textbf{Erro Tipo I} (falso positivo):
rejeitar $H_0$ quando ela, na verdade, era verdadeira.
Ou seja, concluir que existe um efeito que, de fato, não existe.
A chance de cometer esse erro é o próprio nível de significância,
$\alpha$ (geralmente fixado em 0,05).
\item \textbf{Erro Tipo II} (falso negativo):
não rejeitar $H_0$ quando ela, na verdade, era falsa.
Ou seja, deixar passar um efeito que de fato existe.
A chance de cometer esse erro é chamada de $\beta$.
\end{hypeitemize}

\hspace*{1.5em} São erros de natureza oposta,
com causas e consequências distintas,
o que justifica tratá-los como conceitos separados em vez de apenas "o teste errou".

\highlight{
\begin{center}
\resizebox{\linewidth}{!}{
\begin{tabular}{lcc}
\toprule
 & $H_0$ verdadeira & $H_0$ falsa \\
\midrule
Rejeita $H_0$ & Erro Tipo I ($\alpha$) & Decisão correta \\
Não rejeita $H_0$ & Decisão correta & Erro Tipo II ($\beta$) \\
\bottomrule
\end{tabular}
}
\end{center}
}

\hspace*{1.5em} Uma forma consolidada de fixar esses conceitos é pensar num julgamento.
A hipótese nula é "o réu é inocente",
presunção padrão até prova em contrário.
Condenar um inocente corresponde a um Erro Tipo I,
afirmar com convicção algo que não é verdade.
Soltar um culpado corresponde a um Erro Tipo II,
deixar passar algo que de fato aconteceu.
Nenhum sistema de justiça, e nenhum teste estatístico,
consegue eliminar os dois tipos de erro ao mesmo tempo,
existe sempre um trade-off entre eles.

\subtitle{Um exemplo numérico}

\hspace*{1.5em} Considere 1.000 estudos hipotéticos testando um remédio que,
na realidade, não tem efeito nenhum ($H_0$ verdadeira).
Com $\alpha = 0{,}05$,
espera-se que cerca de $1.000 \times 0{,}05 = 50$ desses estudos concluam,
erradamente, que o remédio funciona.

\hspace*{1.5em} Considere agora outros 1.000 estudos testando um remédio que de fato funciona
($H_0$ falsa).
Se o teste utilizado tem poder estatístico de 80\% (ou seja, $\beta = 0{,}20$),
apenas $1.000 \times 0{,}80 = 800$ desses estudos detectam o efeito real,
e $200$ deixam passar,
concluindo erradamente que o remédio não funciona.

\subtitle{Uma ressalva importante}

\hspace*{1.5em} Um equívoco comum é interpretar $\alpha = 0{,}05$ como
"95\% de chance de o efeito ser real".
Não é esse o significado.
$\alpha$ é a chance de errar \textit{dado que} $H_0$ é verdadeira,
valor fixado antes da coleta dos dados,
não calculado a partir deles.
Além disso, $\alpha$ e $\beta$ normalmente estão em tensão,
reduzir o limite para rejeitar $H_0$ (exigir mais evidência, diminuindo $\alpha$)
tende a aumentar a chance de deixar passar um efeito real (aumentando $\beta$),
e o inverso também é verdadeiro.
A forma de reduzir os dois simultaneamente é aumentar o tamanho da amostra.

\begin{infobox}
Em segurança de aeroportos,
um Erro Tipo I corresponde a deter um passageiro sem ameaça alguma (falso alarme).
Um Erro Tipo II corresponde a deixar passar alguém que de fato representa um risco.
Dependendo do contexto,
um erro pode ser consideravelmente mais tolerável que o outro,
o que exige definir previamente qual prioridade faz mais sentido.
\end{infobox}

\originalsource{https://www.youtube.com/watch?v=Hdbbx7DIweQ}

\vspace{0.2cm}

{\color{HypePurple2}Quer se aprofundar mais?}\\
\href{https://errorstatistics.com/wp-content/uploads/2014/08/neyman-pearson-1933_on-the-problem-of-the-most-efficient-tests-of-statistical-hypotheses.pdf}{"On the Problem of the Most Efficient Tests of Statistical Hypotheses"}
(Neyman e Pearson, 1933),
que dá nome ao "Lema de Neyman-Pearson",
usado em testes estatísticos.

}